\documentclass[10pt,a4paper]{amsart}
\usepackage{amsmath,amssymb,amsthm,graphicx}
% Retained minimal preamble for the Stacks Project chapter more-morphisms.tex (More on Morphisms).
% Verbatim excerpt of the chapter preamble: Xy-pic setup (native preamble lines 14-18) and the native
% theorem declarations with the mathematical macros (native preamble lines 154-217). Copied mechanically;
% nothing added, removed or redefined. The wrapper class and the amsmath/amssymb/amsthm/graphicx base
% packages are supplied by the assembler, so the class/externaldocument part of the native preamble is omitted.
% For commutative diagrams we use Xy-pic
\usepackage[all]{xy}

% We use 2cell for 2-commutative diagrams.
\xyoption{2cell}
% Theorem environments.
%
\theoremstyle{plain}
\newtheorem{theorem}[subsection]{Theorem}
\newtheorem{proposition}[subsection]{Proposition}
\newtheorem{lemma}[subsection]{Lemma}

\theoremstyle{definition}
\newtheorem{definition}[subsection]{Definition}
\newtheorem{example}[subsection]{Example}
\newtheorem{exercise}[subsection]{Exercise}
\newtheorem{situation}[subsection]{Situation}

\theoremstyle{remark}
\newtheorem{remark}[subsection]{Remark}
\newtheorem{remarks}[subsection]{Remarks}

\numberwithin{equation}{subsection}

% Macros
%
\def\lim{\mathop{\mathrm{lim}}\nolimits}
\def\colim{\mathop{\mathrm{colim}}\nolimits}
\def\Spec{\mathop{\mathrm{Spec}}}
\def\Hom{\mathop{\mathrm{Hom}}\nolimits}
\def\Ext{\mathop{\mathrm{Ext}}\nolimits}
\def\SheafHom{\mathop{\mathcal{H}\!\mathit{om}}\nolimits}
\def\SheafExt{\mathop{\mathcal{E}\!\mathit{xt}}\nolimits}
\def\Sch{\mathit{Sch}}
\def\Mor{\mathop{\mathrm{Mor}}\nolimits}
\def\Ob{\mathop{\mathrm{Ob}}\nolimits}
\def\Sh{\mathop{\mathit{Sh}}\nolimits}
\def\NL{\mathop{N\!L}\nolimits}
\def\CH{\mathop{\mathrm{CH}}\nolimits}
\def\proetale{{pro\text{-}\acute{e}tale}}
\def\etale{{\acute{e}tale}}
\def\QCoh{\mathit{QCoh}}
\def\Ker{\mathop{\mathrm{Ker}}}
\def\Im{\mathop{\mathrm{Im}}}
\def\Coker{\mathop{\mathrm{Coker}}}
\def\Coim{\mathop{\mathrm{Coim}}}

% Boxtimes
%
\DeclareMathSymbol{\boxtimes}{\mathbin}{AMSa}{"02}

%
% Macros for moduli stacks/spaces
%
\def\QCohstack{\mathcal{QC}\!\mathit{oh}}
\def\Cohstack{\mathcal{C}\!\mathit{oh}}
\def\Spacesstack{\mathcal{S}\!\mathit{paces}}
\def\Quotfunctor{\mathrm{Quot}}
\def\Hilbfunctor{\mathrm{Hilb}}
\def\Curvesstack{\mathcal{C}\!\mathit{urves}}
\def\Polarizedstack{\mathcal{P}\!\mathit{olarized}}
\def\Complexesstack{\mathcal{C}\!\mathit{omplexes}}
% \Pic is the operator that assigns to X its picard group, usage \Pic(X)
% \Picardstack_{X/B} denotes the Picard stack of X over B
% \Picardfunctor_{X/B} denotes the Picard functor of X over B
\def\Pic{\mathop{\mathrm{Pic}}\nolimits}
\def\Picardstack{\mathcal{P}\!\mathit{ic}}
\def\Picardfunctor{\mathrm{Pic}}
\def\Deformationcategory{\mathcal{D}\!\mathit{ef}}
\newenvironment{reference}{}{}
\newenvironment{slogan}{}{}
\newenvironment{history}{}{}
\usepackage{cleveref}
\crefname{originallemma}{original lemma}{original lemmas}
\Crefname{originallemma}{Original lemma}{Original lemmas}
\makeatletter
\newlabel{cohomology-lemma-quasi-separated-cohomology-colimit}{{01FF}{source}{}{sourcepointer-1}{}}
\newlabel{cohomology-lemma-quasi-separated-cohomology-colimit@cref}{{[originallemma][0][]01FF}{[1][0][]source}{}{}{}}
\makeatother
\crefname{originallemma}{original lemma}{original lemmas}
\Crefname{originallemma}{Original lemma}{Original lemmas}
\makeatletter
\newlabel{lemma-globally-generated-vanishing}{{0E7H}{source}{}{sourcepointer-2}{}}
\newlabel{lemma-globally-generated-vanishing@cref}{{[originallemma][0][]0E7H}{[1][0][]source}{}{}{}}
\makeatother
\title{With R\textasciicircum{}1 killed and curve fibres, the structure pushforward is S-flat and commutes with base change}
\author{The Stacks Project Authors}
\date{stacks-project:more-morphisms.tex:local-snapshot-20261011; GNU FDL 1.2 or later}
\begin{document}
\maketitle
\noindent\textit{Editorial note.} Statement span is lines 21192-21206 of sources/stacks-more-morphisms/pinned-original-source.tex; proof span is lines 21208-21232, the complete author proof (23 body lines).\par\medskip
\noindent\textit{Editorial note.} Navigation stubs: 2 cited-but-not-included labels are supplied as non-load-bearing pointers with their pinned tags where the snapshot tag index has them.\par\medskip
\noindent\textit{Editorial note.} extra\_packages is empty on purpose: the retained preamble already loads xy with options and a second declaration clashes in the assembler wrapper.\par\medskip
\noindent\textit{Editorial note.} Text is retained verbatim, including the author's own typos and placeholders.\par\medskip
\noindent\textit{Editorial note.} Third member of the chapter's h1 family (with L16077 propagation and L16097 detection, both collected): this is the flatness-of-H0 theorem the family builds toward. The two author-flagged detail omissions are both one-line standard bookkeeping, every substantive input is named.\par\medskip
\noindent\textit{Allowed standard prerequisites (editorial).}\begin{itemize}
\item Explicit prerequisite boundary. The chapter's lemma globally-generated-vanishing (killing H\textasciicircum{}1 against arbitrary module coefficients); Cohomology, lemma quasi-separated-cohomology-colimit (commutation of cohomology with direct sums on quasi-separated schemes); elementary homological algebra identifying an exact colimit-preserving endofunctor of Mod\_A with tensoring by its value on A.
\end{itemize}\medskip
\ifdefined\hypertarget\hypertarget{sourcepointer-1}{}\fi
\noindent\textit{Source reference.} Original lemma 01FF: sources/stacks-more-morphisms/pinned-original-source.tex, source lines 21232--21232. Cited by the retained statement or proof but not part of the collected spans; supplied only as a navigation pointer. The argument does not depend on its content.\par\medskip
\ifdefined\hypertarget\hypertarget{sourcepointer-2}{}\fi
\noindent\textit{Source reference.} Original lemma 0E7H: sources/stacks-more-morphisms/pinned-original-source.tex, source lines 20971--20980. Cited by the retained statement or proof but not part of the collected spans; supplied only as a navigation pointer. The argument does not depend on its content.\par\medskip
\begin{lemma}
\label{lemma-h1-fibre-zero-h0-flat}
Consider a commutative diagram
$$
\xymatrix{
X \ar[rr]_f \ar[rd] & & Y \ar[ld] \\
& S
}
$$
of morphisms of schemes. Assume $X \to S$
is flat, $f$ is proper, $\dim(X_y) \leq 1$ for $y \in Y$, and
$R^1f_*\mathcal{O}_X = 0$. Then $f_*\mathcal{O}_X$
is $S$-flat and formation of $f_*\mathcal{O}_X$ commutes
with arbitrary base change $S' \to S$.
\end{lemma}

\begin{proof}
We may assume $Y$ and $S$ are affine, say $S = \Spec(A)$.
To show the quasi-coherent $\mathcal{O}_Y$-module
$f_*\mathcal{O}_X$ is flat relative to $S$ it suffices
to show that $H^0(X, \mathcal{O}_X)$ is flat over $A$
(some details omitted).
By Lemma \ref{lemma-globally-generated-vanishing} we have
$H^1(X, \mathcal{O}_X \otimes_A M) = 0$ for every $A$-module $M$.
Since also $\mathcal{O}_X$ is flat over $A$ we deduce the functor
$M \mapsto H^0(X, \mathcal{O}_X \otimes_A M)$ is exact.
Moreover, this functor commutes with direct sums by
Cohomology, Lemma \ref{cohomology-lemma-quasi-separated-cohomology-colimit}.
Then it is an exercise to see that
$H^0(X, \mathcal{O}_X \otimes_A M) = M \otimes_A H^0(X, \mathcal{O}_X)$
functorially in $M$ and this gives the desired flatness.
Finally, if $S' \to S$ is a morphism of affines given by
the ring map $A \to A'$, then in the affine case just discussed
we see that
$$
H^0(X \times_S S', \mathcal{O}_{X \times_S S'}) =
H^0(X, \mathcal{O}_X \otimes_A A') = H^0(X, \mathcal{O}_X) \otimes_A A'
$$
This shows that formation of $f_*\mathcal{O}_X$ commutes
with any base change $S' \to S$. Some details omitted.
\end{proof}

\end{document}
